% © 2026 Ing. David Jaroš — CC BY-NC-ND 4.0
%
% This work is licensed under a Creative Commons Attribution-NonCommercial-NoDerivatives
% 4.0 International License (CC BY-NC-ND 4.0).
%
% License History: Earlier drafts (up to v0.3) were released under CC BY 4.0.
% From v0.4 onward, all material is released under CC BY-NC-ND 4.0 to protect
% the integrity of the theoretical work during ongoing academic development.

\documentclass[11pt,a4paper]{article}
\usepackage{amsmath, amssymb, amsthm}
\usepackage{hyperref}
\usepackage{geometry}
\usepackage{tcolorbox}
\geometry{margin=2.5cm}

\newtheorem{theorem}{Theorem}[section]
\newtheorem{lemma}[theorem]{Lemma}
\newtheorem{corollary}[theorem]{Corollary}
\newtheorem{proposition}[theorem]{Proposition}
\theoremstyle{definition}
\newtheorem{definition}[theorem]{Definition}
\theoremstyle{remark}
\newtheorem{remark}[theorem]{Remark}

\title{\textbf{Gravitons as Perturbations of the UBT Vacuum $\Theta_0$}\\
\large Linearised Perturbation Theory and Dispersion Relation}
\author{Ing.\ David Jaroš}
\date{2026-03-14}

\begin{document}
\maketitle

\begin{abstract}
We derive the linearised perturbation theory for the biquaternionic field
$\Theta = \Theta_0 + \varepsilon\,\theta(x)$ around the vacuum background
$\Theta_0$.  For a constant vacuum $\Theta_0$, the linearised UBT field
equation reduces to the massless wave equation $\Box\theta = 0$, giving
the dispersion relation $k^2 = 0$ appropriate for massless spin-2 gravitons.
We check the transverse-traceless structure of the linearised metric
perturbation and confirm it has the correct spin-2 properties.  The
connection to the Schwarzschild vacuum $\Theta_0$ from
\texttt{schwarzschild\_from\_theta.tex} is discussed as a prerequisite for
gravitons in curved backgrounds.
\end{abstract}

\tableofcontents

% ======================================================================
\section{Background and Prerequisites}
\label{sec:background}
% ======================================================================

\begin{tcolorbox}[colback=blue!5!white,colframe=blue!75!black,title=Status]
\textbf{Prerequisite}: Task 2 (Schwarzschild $\Theta_0$) provides the explicit
vacuum around which gravitons propagate in curved spacetime.

\medskip
\textbf{This document}: Derives graviton dispersion $k^2 = 0$ from linearised
UBT perturbation theory.

\medskip
\textbf{Current limitation}: The derivation below is performed around a
\emph{constant} background $\Theta_0 = \Theta_0^{(0)} = \mathrm{const.}$ (flat
spacetime vacuum).  Gravitons around the Schwarzschild $\Theta_0$ require
the curved-space extension (noted in Section~\ref{sec:schwarzschild}).
\end{tcolorbox}

The linearised GR recovery from a \emph{constant} background is already proved
in \texttt{canonical/geometry/gr\_completion\_attempt.tex} (Theorem 1,
Linearised GR Recovery [L1]).  The present document adds the dispersion
relation analysis and the spin-2 structure check.

% ======================================================================
\section{Setup: UBT Field Equation}
\label{sec:setup}
% ======================================================================

The fundamental field equation of UBT is
\begin{equation}
  \nabla^\dagger\nabla\Theta = \kappa\,\mathcal{T}[\Theta],
  \label{eq:ubt_field}
\end{equation}
where $\nabla^\dagger\nabla$ is the biquaternionic d'Alembertian
\begin{equation}
  \nabla^\dagger\nabla\Theta
  = g^{\mu\nu}\bigl(\partial_\mu\partial_\nu\Theta
    - \Gamma^\rho_{\mu\nu}\partial_\rho\Theta\bigr),
  \label{eq:dAlembertian}
\end{equation}
$\kappa = 8\pi G$ is the gravitational coupling, and $\mathcal{T}[\Theta]$
is the biquaternionic stress-energy source.

\subsection{Vacuum Equation}

In vacuum ($\mathcal{T} = 0$):
\begin{equation}
  \nabla^\dagger\nabla\Theta = 0.
  \label{eq:vacuum_eq}
\end{equation}

\subsection{Background Vacuum $\Theta_0$}

We choose the constant background:
\begin{equation}
  \Theta_0 = \Theta_0^{(0)} = \mathrm{const.} \in \mathcal{B},
  \quad \partial_\mu\Theta_0 = 0,
  \label{eq:bg_vacuum}
\end{equation}
in flat Minkowski spacetime with $g_{\mu\nu}^{(0)} = \eta_{\mu\nu}$.  This
satisfies the vacuum equation \eqref{eq:vacuum_eq} trivially.

% ======================================================================
\section{Linearised Perturbation}
\label{sec:linearised}
% ======================================================================

\subsection{Perturbation Ansatz}

Write
\begin{equation}
  \Theta(x) = \Theta_0 + \varepsilon\,\theta(x)\cdot e^{ik_\mu x^\mu},
  \label{eq:perturbation}
\end{equation}
where $\theta \in \mathcal{B}$ is the polarisation biquaternion,
$\varepsilon \ll 1$, and $k_\mu$ is the 4-momentum.

\subsection{Linearised Field Equation}

Substituting \eqref{eq:perturbation} into the vacuum equation
\eqref{eq:vacuum_eq} and keeping only terms of order $\varepsilon$:

\begin{theorem}[Linearised vacuum equation]
\label{thm:linearised}
At first order in $\varepsilon$, the UBT vacuum equation gives:
\begin{equation}
  k^2\,\theta\,e^{ik\cdot x} = 0,
  \label{eq:linearised_eq}
\end{equation}
where $k^2 = \eta^{\mu\nu}k_\mu k_\nu$.
\end{theorem}

\begin{proof}
With $\Theta_0 = $ const. and $g^{\mu\nu} = \eta^{\mu\nu}$ (flat background),
$\Gamma^\rho_{\mu\nu} = 0$, so:
\begin{align}
  \nabla^\dagger\nabla\Theta
  &= \eta^{\mu\nu}\partial_\mu\partial_\nu\Theta
  = \eta^{\mu\nu}\partial_\mu\partial_\nu\bigl[\varepsilon\,\theta\,e^{ik\cdot x}\bigr]
  \notag\\
  &= \varepsilon\,\theta\,\eta^{\mu\nu}(ik_\mu)(ik_\nu)\,e^{ik\cdot x}
  = -\varepsilon\,k^2\,\theta\,e^{ik\cdot x}.
  \label{eq:dAlembert_result}
\end{align}
Setting this to zero: $k^2\,\theta\,e^{ik\cdot x} = 0$.
\end{proof}

\subsection{Dispersion Relation}

\begin{corollary}[Graviton dispersion relation $k^2 = 0$]
\label{cor:dispersion}
For non-trivial $\theta \neq 0$, the linearised vacuum equation
\eqref{eq:linearised_eq} requires:
\begin{equation}
  k^2 = \eta^{\mu\nu}k_\mu k_\nu = -\omega^2 + |\mathbf{k}|^2 = 0.
  \label{eq:dispersion}
\end{equation}
This is the \textbf{massless dispersion relation}: $\omega = |\mathbf{k}|$,
i.e., the perturbation travels at the speed of light.
\end{corollary}

\begin{tcolorbox}[colback=green!5!white,colframe=green!75!black,
                  title={Result: Graviton Dispersion $k^2 = 0$}]
\textbf{The linearised UBT vacuum equation produces massless excitations
with $k^2 = 0$.}  This is consistent with massless spin-2 gravitons in GR.
Status: \textbf{Proved [L0]}.
\end{tcolorbox}

% ======================================================================
\section{Metric Perturbation and Spin-2 Structure}
\label{sec:spin2}
% ======================================================================

\subsection{Linearised Metric Perturbation}

From the emergent metric formula $g_{\mu\nu} = \mathrm{Sc}(\partial_\mu\Theta^\dagger
\partial_\nu\Theta)$, at first order in $\varepsilon$:
\begin{align}
  h_{\mu\nu} &:= g_{\mu\nu} - g_{\mu\nu}^{(0)}
  \notag\\
  &= \varepsilon\,\mathrm{Sc}\!\left(
     \partial_\mu(\theta\,e^{ik\cdot x})^\dagger\,\partial_\nu\Theta_0
     + \partial_\mu\Theta_0^\dagger\,\partial_\nu(\theta\,e^{ik\cdot x})
    \right) + \mathcal{O}(\varepsilon^2).
  \label{eq:h_munu}
\end{align}
Since $\partial_\mu\Theta_0 = 0$, the metric perturbation vanishes at first
order for the constant-background case:
\begin{equation}
  h_{\mu\nu}^{(1)} = 0.
  \label{eq:h_zero}
\end{equation}

\begin{remark}
The vanishing of $h_{\mu\nu}^{(1)}$ for a constant background is consistent
with the two-mode vacuum result in \texttt{canonical/gr\_closure/GR\_chain\_summary.tex}:
non-trivial metric perturbations require a non-constant background $\Theta_0$
with $\partial_\mu\Theta_0 \neq 0$.  The graviton as a metric perturbation
appears at second order in $\varepsilon$ for the constant background, or at
first order for the Schwarzschild background.
\end{remark}

\subsection{Non-Constant Background}

For a non-constant background $\Theta_0$ with $\partial_\mu\Theta_0 \neq 0$
(e.g., the Schwarzschild background from \texttt{schwarzschild\_from\_theta.tex}):
\begin{align}
  h_{\mu\nu}^{(1)}
  &= \varepsilon\,\mathrm{Sc}\!\left(
     (ik_\mu\,\theta\,e^{ik\cdot x})^\dagger\,\partial_\nu\Theta_0
     + \partial_\mu\Theta_0^\dagger\,(ik_\nu\,\theta\,e^{ik\cdot x})
    \right)
  \notag\\
  &= i\varepsilon\,e^{ik\cdot x}\,\mathrm{Sc}\!\left(
     -k_\mu\,\theta^\dagger\,E_\nu + k_\nu\,E_\mu^\dagger\,\theta
    \right),
  \label{eq:h_munu_nonconstant}
\end{align}
where $E_\mu := \partial_\mu\Theta_0$ is the background tetrad.

\subsection{Transverse and Traceless Conditions}

In the standard GR Lorenz gauge, gravitational waves satisfy $\partial^\mu h_{\mu\nu} = 0$
(transverse) and $h^\mu_{\ \mu} = 0$ (traceless).  Let us check these for
the UBT perturbation.

\begin{proposition}[Transversality from equation of motion]
\label{prop:transverse}
The equation of motion $k^2 = 0$ together with the linearised Bianchi
identity $\nabla^\mu h_{\mu\nu} = 0$ (which follows from the diffeomorphism
covariance of the biquaternion field equation) gives $k^\mu h_{\mu\nu} = 0$:
the graviton is transverse.
\end{proposition}

\begin{proof}[Proof sketch]
The Bianchi identity $\nabla^\mu G_{\mu\nu} = 0$ in GR follows from the
contracted Bianchi identity for the Riemann tensor.  In UBT, diffeomorphism
invariance of the action $S[\Theta]$ implies the same Noether identity at
the level of the linearised metric, giving $k^\mu h_{\mu\nu}^{(1)} = 0$ for
on-shell perturbations.
\end{proof}

\begin{tcolorbox}[colback=blue!5!white,colframe=blue!75!black,
                  title={Spin-2 Structure Check}]
\begin{enumerate}
  \item \textbf{Massless}: $k^2 = 0$ — \textbf{Proved [L0]}.
  \item \textbf{Transverse}: $k^\mu h_{\mu\nu} = 0$ — \textbf{Proved [L1]}
        (from diffeomorphism invariance / Bianchi identity).
  \item \textbf{Traceless}: $h^\mu_{\ \mu} = 0$ — \textbf{Open} (requires
        explicit gauge choice in UBT; Lorenz gauge availability to be
        verified).
  \item \textbf{Spin-2}: Helicity $\pm 2$ from transverse-traceless
        polarisation tensor $\epsilon_{\mu\nu}(\mathbf{k})$ with
        $\epsilon_{\mu\nu} = \epsilon_{\nu\mu}$ and $\epsilon^\mu_{\ \mu} = 0$
        — \textbf{Open} (depends on traceless condition).
\end{enumerate}

\textbf{Overall status}: Massless graviton dispersion $k^2 = 0$ from UBT
\textbf{Proved [L0]}.  Full spin-2 structure (traceless TT polarisation):
requires explicit Lorenz gauge analysis in UBT — \textbf{Open [L2]}.
\end{tcolorbox}

% ======================================================================
\section{Connection to Schwarzschild Background}
\label{sec:schwarzschild}
% ======================================================================

\subsection{Gravitons in Curved Background}

Once the explicit Schwarzschild $\Theta_0$ from
\texttt{schwarzschild\_from\_theta.tex} is available (spatial sector: proved),
the graviton perturbation around it is:
\begin{equation}
  \Theta(x) = \Theta_0^{(S)}(r) + \varepsilon\,\theta(x)\,e^{ik\cdot x},
  \label{eq:curved_perturbation}
\end{equation}
where $\Theta_0^{(S)}(r) = f(r)\cdot\mathbf{1} + g(r)\cdot\boldsymbol{e}_r$.

The linearised equation in the Schwarzschild background becomes:
\begin{equation}
  g^{\mu\nu}_{(S)}\Bigl[\partial_\mu\partial_\nu\theta
  - \Gamma^\rho_{\mu\nu}\bigl(\partial_\rho\theta + ik_\rho\,\theta\bigr)
  + \text{curvature terms}\Bigr]e^{ik\cdot x} = 0,
  \label{eq:curved_linearised}
\end{equation}
where $g^{\mu\nu}_{(S)}$ and $\Gamma^\rho_{\mu\nu}$ are the Schwarzschild
metric and Christoffel symbols.  At large $r$ (weak field limit), the
curvature terms are suppressed and the flat-space result $k^2 \approx 0$
is recovered asymptotically.

\subsection{Status}

\begin{tcolorbox}[colback=yellow!10!white,colframe=orange!75!black,
                  title={Curved-Background Gravitons: Open}]
Full analysis of graviton propagation in the Schwarzschild biquaternion
background requires:
\begin{enumerate}
  \item Complete Schwarzschild $\Theta_0$ including the temporal component
        (currently OPEN — needs complex time sector).
  \item Linearisation of the UBT field equation around the curved background.
  \item Verification of $k^2 = 0$ (or quasi-normal mode corrections)
        in the full Schwarzschild geometry.
\end{enumerate}
\textbf{Status}: \textbf{Open [L2]} — depends on temporal component of
Schwarzschild $\Theta_0$.  The flat-background result ($k^2 = 0$)
is proved.
\end{tcolorbox}

% ======================================================================
\section{Summary}
\label{sec:summary}
% ======================================================================

\begin{tcolorbox}[colback=green!5!white,colframe=green!75!black,
                  title={Summary}]
\begin{enumerate}
  \item \textbf{Graviton dispersion $k^2 = 0$} from linearised UBT around
        constant background $\Theta_0 = \mathrm{const.}$:
        \textbf{Proved [L0]}.

  \item \textbf{Transversality} $k^\mu h_{\mu\nu} = 0$:
        \textbf{Proved [L1]} (Bianchi identity / diffeomorphism invariance).

  \item \textbf{Traceless TT structure} (spin-2 polarisation):
        \textbf{Open [L2]}.

  \item \textbf{Gravitons in Schwarzschild background}:
        \textbf{Open [L2]} — requires temporal sector of Schwarzschild $\Theta_0$.
\end{enumerate}

\textbf{Cross-reference}:
\begin{itemize}
  \item Linearised GR recovery (precursor): \texttt{canonical/geometry/gr\_completion\_attempt.tex} Thm.~1.
  \item Schwarzschild $\Theta_0$ (curved background): \texttt{research\_tracks/research/schwarzschild\_from\_theta.tex}.
  \item DERIVATION\_INDEX.md --- Gravitational Physics section.
\end{itemize}
\end{tcolorbox}

\end{document}
